\item Las ecuaciones 8.43 y 8.44 muestran que $v_{rms} > v_{prom}$ para una colección de partículas de gas, lo cual es cierto siempre que las partículas tengan una distribución de magnitudes de velocidad. Explore esta desigualdad para un gas de dos partículas. Sea $v_1 = v_{prom}$ la rapidez de una partícula y que la otra partícula tenga rapidez $v_2 = (2 - a)v_{prom}$. 
(a) Demuestre que el promedio de estas dos magnitudes de velocidad es $v_{prom}$. 
(b) Demuestre que
$$ v_{rms}^2 = v_{prom}^2(2 - 2a + a^2) $$
(c) Argumente que la ecuación en el inciso (b) demuestra que, en general, $v_{rms} > v_{prom}$. 
(d) ¿Bajo qué condición especial se tendrá $v_{rms} = v_{prom}$ para el gas de dos partículas?
